\documentclass[../main.tex]{subfiles}
\begin{document}

\section{Problema 6 de Pentagrame (Lot 2024)}
\label{problem:6-de-pentagrame}

\href{https://kilonova.ro/problems/2801}{enunț}
$\bullet$
\hyperref[code:6-de-pentagrame]{sursă}

De această problemă nu sînt mîndru. Am scos-o la capăt, dar este lungă și lentă. Soluția oficială folosește doar un arbore de intervale. Eu folosesc și un AIB modificat.

\subsection{Observații teoretice}

Fie $A$, $B$, $C$ și $D$ cadranele de sus-stînga, sus-dreapta, jos-stînga și respectiv jos-dreapta. Să considerăm un $x$ fixat și să vedem cum evoluează partiționările pe măsură ce $y$ crește. Cînd $y = 1$, în mod evident minimul va fi 0 și se va afla în $C$ sau în $D$, iar maximul va fi pozitiv și se va afla în $A$ sau în $B$. Cînd $y = n - 1$ situația se va inversa. Așadar, trebuie să existe:

\begin{itemize}
  \item o primă coordonată $y_1$ unde minimul urcă (trece din $C/D$ în $A/B$);
  \item o primă coordonată $y_2$ unde maximul coboară (trece din $A/B$ în $C/D$).
\end{itemize}

Mai mult decît atît, odată ce minimul urcă, el va rămîne sus, căci $A$ și $B$ continuă să scadă, pe cînd $C$ și $D$ continuă să crească. Similar pentru maxim.

Mai este o observație-cheie care mie mi-a scăpat pînă după încheierea taberei de la Cluj. \emoji{knocked-out-face} Să spunem că la coordonata $y_1$ minimul urcă în cadranul $A$. Atunci pe intervalul dintre $y_1$ și $y_2$ minimul va rămîne în $A$, iar maximul va rămîne în $B$ pînă cînd maximul va coborî. Cele două nu pot face schimb de locuri, căci $A$ continuă să scadă și ar trebui ca $B$ să scadă și mai mult. Or, în clipa în care $B$ nu mai este maxim, înseamnă că maximul a coborît în $C/D$.

Așadar iau naștere patru cazuri, tratate similar:

\begin{enumerate}
  \item $y_1 \leq y_2$, adică minimul urcă primul.

  \begin{enumerate}
    \item Minimul este în $A$, iar maximul este în $B$.
    \item Minimul este în $B$, iar maximul este în $A$.
  \end{enumerate}

  \item $y_1 > y_2$, adică maximul coboară primul.

  \begin{enumerate}
    \item Minimul este în $C$, iar maximul este în $D$.
    \item Minimul este în $D$, iar maximul este în $C$.
  \end{enumerate}
\end{enumerate}

Nu trebuie să uităm despre ce se întîmplă în afara intervalului $[y_1, y_2]$. Aceste cazuri sînt mai ușoare. Pentru $y < y_1$, costul partiționării crește pe măsură ce $y$ scade, căci $A$ și $C$ se dezechilibrează, la fel și $B$ și $D$. Similar, pentru $y > y_2$, costul partiționării crește pe măsură ce $y$ crește. Așadar, putem testa costul partiționării doar în punctele $y_1 - 1$ și $y_2 + 1$. Rămîne partea grea a problemei, testarea între $y_1$ și $y_2$.

\subsection{Structuri de date}

Pentru a afla costul partiționării într-un punct arbitrar, eu am folosit două AIB-uri pentru valorile punctelor înainte, respectiv după $x$. Cu aceeași structură și cu o căutare binară specifică AIB-ului putem căuta pozițiile unde minimul urcă, respectiv unde maximul coboară.

Să considerăm cazul 2.2 din cele patru de mai sus. Care este costul minim al unei partiționări? Avînd în vedere că maximul este în $C$ și minimul în $D$, este suficient să menținem valoarea $C-D$ la fiecare $y$ și să aflăm minimul din $C-D$ pe intervalul $[y_1, y_2]$. Sună a arbore de intervale! Haideți să vedem și cum facem actualizarea.

Vă reamintesc că $C$ și $D$ sînt sume parțiale ale greutăților pentru coordonate $\leq y$, la stînga și respectiv la dreapta lui $x$. Cînd $x$ avansează cu 1, toate punctele care au coordonata egală cu noul $x$ trec din $D$ în $C$. Pentru un punct $P(x_p, y_p, val_p)$, diferența între $C$ și $D$ crește cu $2 \cdot val_p$, iar această creștere afectează cadranele $C$ și $D$ pentru toate partiționările la coordonate $\geq y_p$.

Așadar, avem într-adevăr nevoie de un arbore de intervale cu adăugare pe interval și minim pe interval. Pentru cazul cu minimul în $C$ și maximul în $D$, toate valorile din aint vor fi negative, deci vom avea nevoie de „cea mai puțin negativă” valoare, adică de maxim. Iar pentru minimul și maximul în $A$ și $B$ avem nevoie tot de un arbore de intervale cu sumă pe interval, doar că acum o operație la coordonata $y_p$ va afecta coordonatele $\leq y_p$.

Ca detaliu de implementare, am implementat arborele de intervale și \href{https://kilonova.ro/submissions/376728}{iterativ}. Spre surprinderea mea, sursa este de două ori mai lentă!

\end{document}
